\documentclass{article}

\textwidth=6.5in
\textheight=8.5in
\voffset=-54pt
\oddsidemargin=0in
\evensidemargin=0in
\setlength{\parskip}{8pt}
\setlength{\parindent}{0pt}

\usepackage{amsmath, amssymb}
\usepackage{ifthen}

\usepackage[inline]{enumitem}
\setlist{topsep=0pt, parsep=8pt}

\usepackage[rgb,table]{xcolor}\convertcolorsUfalse
\usepackage{graphicx}

% change hyperref options
\PassOptionsToPackage{hyphens}{url}
\usepackage[bookmarks,colorlinks,linkcolor=black,citecolor=blue,pdfstartview=FitH,urlcolor=blue]{hyperref}
\hypersetup{pdfpagemode=UseNone}
\usepackage{bookmark}

\usepackage{graphicx}
\usepackage{multicol}
\usepackage{framed}
\usepackage{array}

\usepackage{icomma}

\usepackage{marginnote}
\usepackage[colorinlistoftodos]{todonotes}
\renewcommand{\marginpar}{\marginnote}

\usepackage{soul}

\usepackage{pgfplots}
\pgfplotsset{compat=1.18}  
\pgfplotscreateplotcyclelist{mycolor}{black, blue, red, green!70!black, magenta, purple, cyan, orange }  
\pgfplotsset{
  disabledatascaling,
  cycle list=tikzcolor, 
  axis lines=middle,
  every axis plot/.style={no marks, thick},
  every axis/.style={
    enlarge x limits=false,
    enlarge y limits=auto,
    xlabel style={at={(current axis.right of origin)},anchor=west},
    ylabel style={at={(current axis.above origin)},anchor=south},
    % extra x and y ticks for asymptotes
    extra x tick style={grid=major, grid style={dashed,black}},
    extra y tick style={grid=major, grid style={dashed,black}},
    /pgf/number format/1000 sep={},
  },    
  tick label style = {font=\footnotesize, fill=\currentbgcolor, inner sep=0pt},    
  xticklabel shift=1pt,    
  scaled ticks=false,
  scaled y ticks=false,
  unbounded coords=jump,
  samples=101,
  minor tick length=0,
  scatter/classes={
    open={mark=*,draw=black,fill=white, mark size=2, thin},
    closed={mark=*,draw=black,fill=black, mark size=2, thin}
  },
  width=3in,
}
\pgfplotsset{open/.style={only marks, mark=*, fill=white, thin, mark size=2}}
\pgfplotsset{closed/.style={only marks, mark=*, draw=black, fill=black,
    thin, mark size=2}}
\pgfplotsset{labeled/.style={nodes near coords, only marks, mark=*, draw=black, fill=black, thin, mark size=2, point meta=explicit symbolic}}

\usepackage{tikz}
\usetikzlibrary{
  matrix,%
  % enables the snake paths
  decorations.pathmorphing,%
  % enables bigger arrows
  decorations.markings,%
  decorations.pathreplacing,%
  decorations.text,%
  calc,%
  patterns,%
  shapes.geometric,%
  arrows,
  decorations.shapes,
  positioning,
  plotmarks,
  intersections
}
\tikzset{>=latex} % sets default arrow type in tikz
\tikzset{font=\normalsize}
\tikzset{mark size=1.5pt, mark options=thin}
\tikzset{pin distance=4pt,
  every pin edge/.style={<-, thin, shorten <= -2pt}}

\interfootnotelinepenalty=10000
\renewcommand{\thefootnote}{(\arabic{footnote})}

\usepackage{multirow, bigdelim}

% change to \usepackage[sols]{handout} to see solutions
\usepackage[sols, grading, notes]{handout}

\newcommand\iscurrentenumi[1]{%
  \ifthenelse{\equal{#1}{\theenumi}}{}{#1}}
\setenumerate[1]{ref=\#\arabic*}
\setenumerate[2]{ref=\protect\iscurrentenumi{\theenumi}(\alph*)}
\newcommand\enumiiprefix[2]{%
  \ifthenelse{{\equal{#1}{\theenumi}}\AND{\equal{#2}{(\alph{enumii})}}}{}{}%
  \ifthenelse{{\equal{#1}{\theenumi}}\AND{\NOT{\equal{#2}{(\alph{enumii})}}}}{#2}{}%
  \ifthenelse{\equal{#1}{\theenumi}}{}{#1#2}%
}
\setenumerate[3]{ref=\protect\enumiiprefix{\theenumi}{(\alph{enumii})}\roman*}

\makeatletter

% underline links               
\let\@href=\href
\renewcommand{\href}[2]{\@href{#1}{\ul{#2}}}

\makeatother

\definecolor{purple}{rgb}{0.5,0,1.0}

\newcommand{\doedfinity}[2][\newpage]{
  \begin{problem}[#1]
    Do the Edfinity assignment ``Problem Set #2''.

    We encourage you to write down any scratch work here, as well as
    things you want your future self to remember. Nothing you write
    here will be graded; it's just for you!
  \end{problem}
}

\newcommand{\psrnewpage}{\newpage}
\newcommand{\psreflection}[2][]{

  \ifthenelse{\not\equal{#1}{nocite}}{
    \begin{enumerate}[resume]
    \item
      \begin{problem}[1in]
        Please cite any help you got on this problem set:
        \begin{itemize}[nosep]
        \item If you worked with other students, please list their names here.
        \item If you used resources other than those on our Canvas site, please list them here.
        \end{itemize}
      \end{problem}

    \end{enumerate}
  }

  \probonly{%
    #2

    \psrnewpage

    \emph{Reflection space.} It's a good habit to take a few minutes at the end of each pset to reflect on what you learned and jot down notes to your future self. Here are a few reflection questions to get you started. Nothing you write here will be graded; it's just for you!
    \begin{itemize}
    \item Take a few minutes to look back at the problems. How could you explain to someone else how to do each problem? What was your strategy, and how did you know to use that strategy? If there are problems where you don't feel able to answer this, write down which ones they are. Come back to them in a few days when we post the homework solutions!

      \vspace{1.5in}

    \item Take a look back at the learning objectives on today's worksheet; how are you feeling about each one? If there are ones you know you need more practice with, write those down here.

      \vspace{1.5in}

    \item What did you find most challenging about this material? Do you have any lingering questions about it? If so, write them down here so you don't forget them. At some point, stop by office hours or MQC to ask!

      \vfill
    \end{itemize}      
  }
}

\usepackage{fancyhdr}
\lhead{\textcolor{white}{This content is protected and may not be shared, uploaded, or distributed.}}
\rhead{}
\rfoot{\vspace*{\fill}\footnotesize\textcopyright{} \emph{Harvard Math Ma}}
\renewcommand{\headrulewidth}{0pt}
\pagestyle{fancy}
\begin{document}

\begin{center}
  \large\textsc{Problem Set 36 - Comparing Exponentials, Logarithms, and Power Functions}
\end{center}

\begin{enumerate}
\item \doedfinity{36}

\item  

  \begin{enumerate}
  \item \label{growth-rate-limits-1}

    \begin{grading}
      2 points per part: 0.5 for the limit, 1.5 for the explanation. If they use notation like $\infty^2$ or $e^\infty$ as part of their justification, please deduct 0.5. (If they've written something like $\infty^2$ as scratch work but then have a complete justification that uses correct notation, don't deduct.)
    \end{grading}

    Evaluate the following limits. Please justify carefully, without using mathematically meaningless expressions like $e^\infty$ and $\infty^2$.

    \begin{enumerate}
    \item
      \begin{problem}[\vfill]
        $\displaystyle \lim_{x \to \infty} \dfrac{x^2}{2x^2 + e^x}$
      \end{problem}

      \begin{solution}
      As $x \to \infty$, the dominant term in the denominator is
      $e^x$.
        \begin{align*}
          \lim_{x \to \infty} \frac{x^2}{2x^2 + e^x} &= 
          \lim_{x \to \infty} \frac{x^2}{e^x} \\
          &= 0 \text{ since $x^2 \ll e^x$}
        \end{align*}
      \end{solution}

    \item
      \begin{problem}[\vfill]
        $\displaystyle \lim_{x \to -\infty} \dfrac{x^2}{2x^2 + e^x}$
      \end{problem}

      \begin{solution}
      As $x\to -\infty$, the dominant term in the denominator
      is $2x^2$. 
        \begin{align*}
          \lim_{x \to -\infty} \frac{x^2}{2x^2 + e^x} &= 
          \lim_{x \to -\infty} \frac{x^2}{2x^2} \\
          &= \lim_{x \to -\infty} \frac{1}{2} \\
          &= \frac{1}{2}
        \end{align*}
      \end{solution}

    \item
      \begin{problem}[\vfill]
        $\displaystyle \lim_{x \to \infty} \dfrac{3 e^{2x} - x - 3}{e^{2x} + 1}$
      \end{problem}

      \begin{solution}
        As $x to \infty$ the dominant terms in the numerator
        and denominator are $3e^{2x}$ and $e^{2x}$ respectively.
        \begin{align*}
          \lim_{x \to \infty} \frac{3 e^{2x} - x - 3}{e^{2x} + 1} &=
          \lim_{x \to \infty} \frac{3 e^{2x}}{e^{2x}} \\
          &= \lim_{x \to \infty} 3 \\
          &= 3
        \end{align*}
      \end{solution}

    \item
      \begin{problem}[\vfill]
        $\displaystyle \lim_{x \to -\infty} \dfrac{3 e^{2x} - x - 3}{e^{2x} + 1}$
      \end{problem}

      \begin{solution}
        As $x\to-\infty$, the dominant terms in the numerator
        and denominator are $-x$ and $1$ respectively.
        \begin{align*}
          \lim_{x \to -\infty} \frac{3 e^{2x} - x - 3}{e^{2x} + 1}
          &= \lim_{x \to -\infty} \frac{-x}{1} \\
          &= \infty
        \end{align*}
      \end{solution}

    \end{enumerate}

  \item

    \begin{grading}
      1 point per part: 0.25 for choice, 0.75 for explanation
    \end{grading}

    In each part, you're given a function; without using technology, pick the picture (choices A - D on the next page) that shows its graph. Explain your choice. (How is (a) helpful?)

    \begin{enumerate}
    \item
      \begin{problem}[\nstretch{0.75}]
        $f(x) = \dfrac{x^2}{2x^2 + e^x}$
      \end{problem}
      \begin{solution}
        We found in part (a) that $f(x)$ approaches $0$
        as $x \to \infty$ and $f(x)$ approaches $\frac{1}{2}$
        as $x \to -\infty$. Of the choices, (C) is the best match.
      \end{solution}

    \item
      \begin{problem}[\vfill]
        $g(x) = \dfrac{3 e^{2x} - x - 3}{e^{2x} + 1}$
      \end{problem}
      \begin{solution}
        We found in part (a) that $g(x)$ approaches $3$
        as $x \to \infty$ and $g(x)$ approaches $\infty$
        as $x \to -\infty$. Of the choices, (A) is the best match.
      \end{solution}

    \end{enumerate}

    Choices:
    \begin{center}
      \begin{tabular}{c c c c c c c}
        \begin{tikzpicture}
          \begin{axis}[xtick=\empty, ytick=\empty, width=1.8in]
            \addplot+[domain=-8:7] {(3*e^(2*x) - x - 3) / (1 + e^(2*x))}; 
          \end{axis}
        \end{tikzpicture}
        &&
        \begin{tikzpicture}
          \begin{axis}[xtick=\empty, ytick=\empty, width=1.8in]
            \addplot+[domain=-4:0] {3/(1+x^2)-3};
            \addplot+[domain=0:3] {x^2};
          \end{axis}
        \end{tikzpicture}        
        &&
        \begin{tikzpicture}
          \begin{axis}[xtick=\empty, ytick=\empty, width=1.8in]
            \addplot+[domain=-4:9] {x^2/ (2*x^2 + e^(x))}; 
          \end{axis}
        \end{tikzpicture}
        &&
        \begin{tikzpicture}
          \begin{axis}[xtick=\empty, ytick=\empty, width=1.8in]
            \addplot+[domain=-6:10] {(e^(-3 + x) + e^(2*(-3 + x))-e^(-3)-e^(-6))/(1 + e^(2*(-3 + x)))}; 
          \end{axis}
        \end{tikzpicture} \\
        (A) && (B) && (C) && (D)
      \end{tabular} 
    \end{center}
  \end{enumerate}

\item Evaluate the following limits. Again, justify your answers carefully, using correct notation.

  \begin{grading}
    2 points per part, same breakdown as in \ref{growth-rate-limits-1}
  \end{grading}

  \begin{enumerate}
  \item
    \begin{problem}[\vfill]
      $\displaystyle \lim_{x \to \infty} \frac{\ln x + \sqrt[10]{x}}{x^{0.1}}$
    \end{problem}

    \begin{solution}
      As $x \to \infty$, the dominant term in the numerator is 
      $\sqrt[10]{x}$.
      \begin{align*}
        \lim_{x \to \infty} \frac{\ln x + \sqrt[10]{x}}{x^{0.1}}
        &= \lim_{x \to \infty} \frac{\sqrt[10]{x}}{x^{0.1}} \\
        &= \lim_{x \to \infty} \frac{x^{0.1}}{x^{0.1}} \\
        &= 1
      \end{align*}
    \end{solution}

  \item
    \begin{problem}[\vfill]
      $\displaystyle \lim_{x \to \infty} \frac{\ln x + 0.99^x}{\sqrt[100]{x}}$ 
    \end{problem}

    \begin{solution}
      As $x\to \infty$, the dominant term in the numerator
      is $\ln x$.
      \begin{align*}
        \lim_{x \to \infty} \frac{\ln x + 0.99^x}{\sqrt[100]{x}}
        &= \lim_{x \to \infty} \frac{\ln x}{x^{0.01}} \\
        &= 0 \text{ since $\ln x \ll x^{0.01}$}
      \end{align*}
    \end{solution}

  \item
    \begin{problem}[\vfill\newpage]
      $\displaystyle \lim_{x \to \infty} \frac{\ln x + 1.01^x}{x^5}$
    \end{problem}

    \begin{solution}
      As $x\to\infty$, the dominant term in the numerator
      is $1.01^x$.
      \begin{align*}
        \lim_{x \to \infty} \frac{\ln x + 1.01^x}{x^5}
        &= \lim_{x \to \infty} \frac{{1.01}^x}{x^5} \\
        &= \infty \text{ since $1.01^x \gg x^5$}
      \end{align*}
    \end{solution}

  \end{enumerate}

\item 

  \begin{grading}
    7 points:
    \begin{itemize}[nosep]
    \item 2 points for correctly finding critical points of $f$
    \item 1 for correctly evaluating $\displaystyle \lim_{x \to \infty} f(x)$
    \item 2 for correctly $\displaystyle \lim_{x \to -\infty} f(x)$ and justifying their limit. (If they don't say anything about the sign of the denominator, please deduct 1.)
    \item 2 for comparing these limits with the value of $f$ the critical points and drawing correct conclusions
    \end{itemize}
    If they don't check any limits but correctly use the First or Second Derivative Test to classify 0 as a local minimum and 4 as a local maximum, give 3.5 points out of 7. 
  \end{grading}

  \begin{problem}[\newpage]
    Does $f(x) = \dfrac{x^4}{e^x}$ have a global maximum on $(-\infty, \infty)$? If so, where? If not, why not? How about a global minimum? Justify your answers carefully and completely. 
  \end{problem}

  \begin{solution}
    First we need to find the critical points.
    We can use the Quotient Rule
    to find $f'(x)$:
    \begin{align*} 
    f'(x) &= \frac{(4x^3)(e^x) - (e^x)(x^4)}{(e^x)^2} \\
    &= \frac{4x^3 - x^4}{e^x} \\
    &= \frac{x^3(4-x)}{e^x}
    \end{align*}
    $f'(x)$ is
    never undefined, and $f'(x)=0$ when $x = 0, 4$.  
    So the critical
    points are $0$ and $4$. 
    We can determine the global extrema by
    comparing the value of $f(x)$
    at these critical points with $\lim\limits_{x\to \infty} f(x)$
    and $\lim\limits_{x \to -\infty} f(x)$.
    \begin{itemize}
        \item $f(0)=0$
        \item $f(3)=\dfrac{81}{e^3}$
        \item $\lim\limits_{x\to \infty} f(x) = 0$, since $x^4 \ll e^x$
        \item $\lim\limits_{x \to -\infty} f(x) = \infty$ since $x^4$
        approaches $\infty$ and $e^{x}$ is positive and approaching $0$.
    \end{itemize}
    Based on these values, we can see that there is a global minimum
    at $(0, 0)$, and there is no global maximum.
  \end{solution}

\item  \begin{enumerate}
  \item

    \begin{grading}
      3 points: 1 for a reasonable description of odd function, 1 for reasonable description of polynomial with odd degree, 1 for an example of an odd function that isn't a polynomial
    \end{grading}

    \begin{problem}[\vfill]
      A classmate says to you, ``I'm confused about the difference between an odd function and a polynomial with odd degree.'' Explain what the terms \emph{odd function} and \emph{polynomial with odd degree} mean, and give an example of an odd function that isn't a polynomial. 
    \end{problem}

    \begin{solution}
      An odd function is a function whose graph has symmetry
      around the origin. That is, $f(-x) = -f(x)$. A polynomial with
      odd degree is a polynomial like $x^3 + 1$, where the highest
      exponent of the variable is odd. One example
      of an odd function that is not a polynomial is $f(x)=\dfrac{1}{x}$.
    \end{solution}

  \end{enumerate}

  Decide whether each statement below is true or false. If the statement is true, explain why. If it's false, give a counterexample (an example demonstrating that the statement is false); your counterexample can be a formula or graph. 

  \begin{grading}
    1.5 points per part: 0.5 for correctly saying whether it's true or false, 1 for explanation / counterexample. 
  \end{grading} 

  \begin{enumerate}[resume]

  \item
    \begin{problem}[\vfill]
      If a polynomial $p(x)$ is invertible, then it must have odd degree. 
    \end{problem}

    \begin{solution}
      This is \fbox{true}. If a polynomial function $p(x)$ is invertible,
      then it must pass the vertical line test, so it must be always
      increasing or always decreasing. This implies that the polynomial
      must have odd degree, since a polynomial with even degree
      has at least one turning point.
    \end{solution}

  \item
    \begin{problem}[\vfill]
      If a polynomial $p(x)$ has even degree, then it must be an even function. 
    \end{problem}

    \begin{solution}
      This is \fbox{false}. A counterexample is $f(x)=x^2+x$.
    \end{solution}

  \item
    \begin{problem}[\vfill\newpage]
      If a polynomial $p(x)$ is an even function, then it must have even degree.
    \end{problem}

    \begin{solution}
      This is \fbox{true}. If a polynomial function satisfies
      $p(-x) = p(x)$, then there cannot be \emph{any} odd powers
      of $x$ when $p(x)$ is expressed in standard form.
    \end{solution}

  \end{enumerate}

\item 

  \begin{grading}
    3 points, full credit for any thoughtful response
  \end{grading}

  The final exam is on Monday 12/18. Please read  \href{https://canvas.harvard.edu/courses/125618/pages/exams}{the information on Canvas about the exam}.

  \begin{enumerate}
  \item

    \begin{problem}[\newpage]
      We will have several review sessions. Take a look at the \href{https://canvas.harvard.edu/courses/125618/pages/office-hours-and-review-sessions}{schedule}; review sessions are highlighted in yellow.

      Don't try to attend every single session! Instead, pick ones on topics that you've struggled with. For example, if you had trouble with Skills Check 2, then it makes sense to go to the session on exponentials \& logarithms; if you had trouble with the optimization problem on Exam 2, check out the sessions on optimization.

      We're offering a proctored practice exam on Monday 12/11; this is a chance to take a practice exam in a test-like setting, and it's a good way to see what topics you need to work on more. If you'd like to take the proctored practice exam, please \href{https://forms.gle/y4QdH2kqt748Dfua9}{RSVP} by Friday 12/8 so that we can be sure to have enough copies.

      Which sessions do you plan to attend? 
    \end{problem}

  \item
    \begin{problem}[\vfill]
      Look back at the \href{https://docs.google.com/presentation/d/e/2PACX-1vTdTTLeI90WdBVbxU-72oaJ5_7bJSrwMDNSPTPxYl3FZw1nwCkcD-csVSR8SFRkGNGszXwssfBJEd2D/pub?start=false&loop=false&delayms=60000}{slides from workshop on using reading period effectively}. Using the advice there, work backwards from our exam date (Monday 12/18) to make a schedule for your final exam studying.

      We've posted 4 practice exams, and we strongly encourage you to do at least 3 of them. 
    \end{problem}
  \end{enumerate}

\end{enumerate}
\psreflection{}
\end{document}
